% Generated by scripts/08_extract_bibliography.py.
\section*{Oracle-Bone Collections Cited and Their Abbreviations}
\addcontentsline{toc}{section}{Oracle-Bone Collections Cited and Their Abbreviations}

\begin{description}
  \item[Cheng-chai] Sun Hai-po 係 海 波 , Ch’eng-chai Yin-hsii wen-tzu 誠 說 股 虎 多 了 字. Peking. 1940. \hfill\textit{[CJK: 小屯第二本 殷壶文字 甲編考釋]}
  \item[Ghi-ch] eng 1 顏一萍 最 —-j\#. Chia-ku chi-ch’eng Wy BAK . Vol. 1. [Taipei], 1973: |
  \item[Chi-chi] Fang Fa-han 75 1 tk \{Frank H. Chalfant] and Po Jui-hua Si i€ [Roswell | on S. Britton]. Chia-ku pu-tz’u ch’i-chi 4p b AY LR Seven Collections of Inscribed | Oracle Bone\}. New York, 1938. Reprint, Taipei, 1966. \hfill\textit{[CJK: 小屯第二本 殷壶文字 甲編考釋]}
  \item[甲編] Tung Tso-pin ffi ff. Hsiao-tun ti-erh-pen : Yin-hsii wen-tzu: chia-pien 小 屯 第 A: By uke 0g? Il Ha. Nanking, 1948. Reprint, Taipei, 1977. \hfill\textit{[CJK: 小屯第二本 殷壶文字 甲編考釋]}
  \item[甲編] 99- 訪 沒 Chu Wan-li i ag. Asiao-Vun t-erh-pen: Yin-hsii wen-tzu: chia-pien k'ao-shih \~{} LA AR Bea Ce 4 BE. Taipei, 1961. 229 \hfill\textit{[CJK: 小屯第二本 殷壶文字 甲編考釋]}
  \item[Chia-tu—] T u-pan \textbackslash\{\}ili (Plates 1-50 at the rear of 甲編 k’ao-shih).
  \item[Chien-shou] Chi Fo-t’o iti (@8 §E . Chien-shou-tang so-ts’ang Yin-hsii wen-tzu (ih 2h 1 NV aX BC LF Shanghai, 1917. Reprint, [Taipei], n.d.
  \item[Givien-pien] Lo Chen-yu fie i. Vin-hsié shu-ch’i chien-pien Bx ite HR wii. N.p., 1913. | Reprint, [Shanghai], 1932; [Taipei], n.d.
  \item[Chih-hsii] Li Ya-nung \#ti lft. Yin-ch't chih-yi hsii-pien (i 2. He (ik Had. Peking, 1950. Chikevi Li Tan-ch’iu 4H. Yin-ch’t chih-yi eR Hk. Peking, 1941.
  \item[Chin-chang] Fang Fa-han [Frank H. Chalfant] and Po Jui-hua [Roswell S. Britton]. Chin- chang so-ts’ang chia-ku pu-tz’u 4 StF we HT b BE (The Hopkins Collection of Inscribed Oracle Bone\}. New York, 1939. Reprint, Taipei, 1966. \hfill\textit{[CJK: 收集 殷文字合]}
  \item[京津] Hu Hou-hsiian 明 厚 宣 , Chan-hou 京津 hsin-huo chia-ku-chi We RCE MEH! Shanghai, 1954. \hfill\textit{[CJK: 收集 殷文字合]}
  \item[Ghing-hua] Lo Chen-yii. Yin-hsii shu-ch’t ching-hua Wit BRL He. N.p., 1914. tet Cho-ts'un Tseng Yi-kung \$4 \#4 . Chia-ku cho-ts’un \textbackslash\{\}fl tp 3\% ff. N.p-, 1939. Chuizho Kuo Jo-vii 7 \#8 ; Tseng Yi-kung; and Li Hstieh-ch'in \#4! 4). Yin-hsti wen-tzu \{ chui-ho Bye ike 7:7. Peking, 1955.
  \item[Chui-hsin] 顏一萍. Chia-ku chui-ho hsin-pien 甲骨 弧 合 新 編 . Taipei, 1975. oat Fu-chia Shang Ch’eng-tso pq Kile. Fu-shth so-ts’ang chia-ku wen-tzu wi KOT xP a oc - Nanking, 1933. \hfill\textit{[CJK: 漢城大學所藏大胛骨刻辭考釋]}
  \item[Fu-yin] Wang Hsiang |: \$8. Fu-shih Yin-ch’t cheng-wen (83 84 tk 4c. Tientsin, 1925. \hfill\textit{[CJK: 漢城大學所藏大胛骨刻辭考釋]}
  \item[Hai-wai] Jao Tsung-vi 饒宗頤 .“Hai-wai chia-ku lu-yi 海外 甲骨 錄 貴 . Journal of Oriental Studies 4.1-2 (1957/58), pp. 1-22. : Han-ch'eng Tung Tso-pin. “Han-ch’eng ta-hsiieh so-ts’ang ta-chia-ku k’e-tz’u k’ao-shih \hfill\textit{[CJK: 漢城大學所藏大胛骨刻辭考釋]}
  \item[TER] ACE AE AINE SEAR.” BIHP 28.2 (1957), pp. 825-840. 1 Ho-pien Tseng Yi-kung. Chia-ku chui-ho-pien 甲骨 組 合 編 . N.p., 1950.
  \item[Hou-pien] Lo Chen-yii. Yin-hsii shu-ch’t hou-pien Ft itt 書 契 後 編 . N.p., 1916. Reprint, [Taipei], n.d. Le it \} i Hsia-men Hu Hou-hsiian. “‘Hsia-men ta-hsiieh so-ts’ang chia-ku wen-tzu 說 門人 學 所 藏 甲 , | i HCE In Chia-ku-hsiieh Shang-shih lun-tsung ch’u-chi if 44/8 pg lt ite REA | | i Ch’eng-tu, 1944. 還 |
  \item[Asii-pien] Lo Chen-yti. Yin-hsii shu-ch’i hsii-pien JRE i itd N.p.. 1933. Reprint, , ! [Taipei], n.d. ry |
  \item[Hsii-pu] | Hu Hou-hsiian. Chia-ku hsii-pu 貼 全 續 補 [unpublished collection. title uncertain]. | | ; Hsii-ts-un Hu Hou-hsiian. Chia-ku hsii-ts’un Hlth ete. Shanghai, 1955. 7 人文 Kaizuka Shigeki 14g SBT . Kydto daigaku jimbun hagaku kenkyiijo 20 kokotsu monji
  \item[I:] FAD AAR A SCRE SE STE TT 0\% © 2 vols. Kyoto, 1959. Shakubun Rc. 圖 | Kyoto, 1960. Sakuin 索引 . Kyoto, 1968. i | | Kikké Hayashi Taisuke 林 Agi\}. Kikkd jitkotsu monji PMP Ye see N.p.. 1921. Reprint, | Taipei, 1970. “| \hfill\textit{[CJK: 中國文字 睦之輔氏藏甲骨文字。 東方學報]}
  \item[Keu-fang] Fang Fa-lien Jj\}: [Frank H. Chalfant] and Po Jui-hua [Roswell S. Britton]. ; ; Kou-Fang erh-shih-ts’ang chia-ku pu-tz Wi Sj AEM Pb RE (The Couling-Chalfant | i Collection of Inscribed Oracle Bone). Shanghai, 1935. Reprint, Taipei, 1966. 7 | : Ling-shih Ch'en Pang-huai [dt . Chia-ku-wen ling-shih ith scat He | Tientsin, 1959. ; | Reprint, Tokyo, 1970. ! | \hfill\textit{[CJK: 甲骨學史論叢三集]}
  \item[Liu-lu] Hu Hou-hsiian, “Chia-ku liu-lu H! 4\} A... In Chia-ku-hsiieh Shang-shth lun-ts’ung f san-cha "lay FY ng We it 22 HE. Ch’eng-tu, 1945. Ht 全 Menzies sii Chin-hsiung. The Menzies Collection of Shang Dynasty Oracle Bones. Volume 1: , A Catalogue. Toronto [1972]; Volume IT: The Text. Toronto (1977).* ) | : i Ming-hou Ming Yi-shih 84\% 4: [James M. Menzies]. Yin-Asii pu-tz’u hou-pien IB Mi b BY A AE. 4 | a Edited by Hsii Chin-hsiung 許 進 雄 . 2 vols. Taipei, 1972. an it Na-erh 顏一萍. Mei-kuo Na-erh-sen mei-shu-kuan ts’ang chia-ku pu-tz'u k’ao-shih 美 國 1 由 230 AA GR GR SE Or HT ake HE Pb tA Taipei, 1973. (Originally published in Chung-kuo | | ro! 同 wen-lzu 中 國 文 字 22--25, 29 [1966- 1968].) \~{} | bo Nan-pet Hu Hou-hsiian. Chan-hou nan-pet so-chien chia-ku-lu ®\& ff \textbackslash\{\}44 IEE SE BY ER. Peking | my and Shanghai, 1951. iyo 1 ‘. 寧滬 Hu Hou-hsiian. Chan-how 寧滬 hsin-huo chia-ku-chi ‘Wy (Wt ME LEP AS . | | i Peking, 1951. 下, | th Ogawa td Michiharu ff if iii i . “Ko Ogawa Chikanosuke shi 26 kokotsu monji 故 小 川 ! fh eZ SCH TS CF.” 7227 gakuhd 東方 學 報 37 (1966), pp. 249-263. | be Pa-li Jao Tsung-yi. Pa-li so-chien chia-ku-lu E 2) FAY Bk. Hongkong, 1956. wd \hfill\textit{[CJK: 甲骨學史論叢三集 中國文字]}
  \item[Pai-ken] Ming Yi-shih [James M. Menzies]. Pai-ken-shih chiu-ts’ang chia-ku wen-tzu ff Hi EC vi Mm ce (The Paul D. Bergen Collection: Chinese Oracle Bone Characters). : 1 Ch’i-ta chi-k’an 齊 人 季刊 6and 7 (1935). | | | 7 拼編 Chang Ping-ch’iian 48 9¢ RE. 小屯 ti-erh-pen: Yin-hsii wen-tzu; ping-pien 小 子 L : : 股 虛 交 字 : 丙 編 .Talpei. Vol. 1, pt. 1 (1957); pt. 2 (1939). Vol. 2, pt. ! 1 (1962); pt. 2 (1965). Vol. 3, pt. 1 (1967); pt. 2 (1972). j te Pu-tzu Jung Keng 容 護 and Ch’é Jun-min \#2 i . Yin-ch’i pu-tz’u BY | AH . Peiping, | | 1933. Reprint, Taipei, 1970. | 4 ROM Royal Ontario Museum. [Used to refer to oracle bones in that collection; also | ! used by some as an abbreviation equivalent to Menzies. | | \hfill\textit{[CJK: 中國文字]}
  \item[Shan-chai] Liu T’i-chih 劉 體 智 . Shan-chai chiu-ts’ang Pa-pen BIR BBL. Unpublished. | | 1 Shih-to Kuo Jo-yii. Yin-ch’i shih-to BY 32 18 \$8. Vol. 1. Shanghai, 1951. Vol. 2. Peking, 1953. 了 i T’ai-ta tr Tung Tso-pin. ‘‘Tai-wan ta-hsiieh so-ts’ang chia-ku wen-tzu fu-k’ao-shih 4 \}f¥ | i 大 學 所 藏 甲骨 文 字 附 考 釋 .” Kuo-li T’ai-wan ta-hsiieh k’ao-ku jen-let hsiieh-k’an | | i eMAB eH AMF + (1953), pp. 22-46. | e T’ai-ta 2. Tung Tso-pin and Chin Hsiang-heng 金 祥 恆 .““Pen-hsi so-ts’ang chia-ku wen- “| P *The second volume was received as Sources was in galleys; unless otherwise indicated, all references to Menzies in this 一 book are to Volume I. ; J | AA 4 7 ww tau AR Pap sce Kuo-li T’ai-wan ta-hstieh k'ao-ku jen-lei hsiieh-k'an 17-18 we (1961), pp. 71-84. | \hfill\textit{[CJK: 國立臺灣大學考古人類學刊 甲骨學史論叢三集]}
  \item[T\textquotesingle{}ieh-hsin] 顏一萍. Tieh-yin ts’ang-kuei hyin-pien WUE Ga Hf Bi. Taipei, 1975. |
  \item[Tieh-shih] Li Tan-ch’iu. Tieh-yiin gg ling-shih BUS RW AAG . Shanghai, 1939.
  \item[Tieh-yi] Yeh Yii-sen 葉 玉 森 . Tieh-viin 生 gg shih-yi BREE BGI N.p., 1925. Techy Lo Chen-yit. Tieh-yiin 全 \~{} chih-yit BUR BZ HE. N.p., 1915. Reprint, | Hongkong, 1972. |
  \item[Tieh-yiin] Liu E898 . Tieh-yiin ts’ang-kuei Bee wi 1. N.p., 1903. Reprint, [Taipei], 1959. |
  \item[Tien-jang] T'ang Lan 店 蘭 . Tten-jang-ke chia-ku wen-ts’un K EMR) Hy sc FE Peiping, 1939. bw Ts‘ui-pien Kuo Mo-jo \$0:4 \# . Yin-ch’i ts’ui-pien 1 22\% \# . Tokyo, 1937. Rev. ed., Peking, 1959. Reprint. Taipei, 1971.
  \item[Ts’un-chen] Kuan Pai-yi \$\$ ¢143 . Yin-hsii wen-tzu ts’un-chen EME ACF FEI . Vols. 1-5, 8. |
  \item[K’ai-feng] 1931-1938. | a T'ung-tsuan Kuo Mo-jo. Pu-tz’u Pung-tsuan \} Ysii \&\% . Tokyo, 1933. |
  \item[CS] 周洪祥. Oracle Bone Collections in the United States. Occasional Papers | in Archaeology. no. 10. Berkeley and Los Angeles, 1976. [USB refers to a bone fragment in this collection, USS refers to a shell fragment. ]
  \item[Wai-pien] Tung Tso-pin. Yin-hsii wen-tzu wai-pien BB ike L7H. Taipei, 1956.
  \item[Wen-lu] Sun Hai-po. Chia-ku wen-lu A\} 3\%. Honan, 1938. Reprint, Taipei, 1971. !
  \item[White] sii Chin-hsiung. Oracle Bones from the White and Other Collections [Forthcoming. ] | loa Yeh-ch’u Huang Chin i . Yeh-chung p’ien-yii ch’u-chi ¥b rh Jt 4 . Peiping, 1935. [The oracle-bone rubbings are reprinted in Yeh-chung p’ien-yii. Taipei, 1972.]
  \item[Yeh-san] Huang Chiin. Yeh-chung p’ien-yii san-chi \#5 +P fy = 9. Peiping, 1939. [The oracle-bone rubbings are reprinted in Yeh-chung p’ien-yii. Taipei, 1972.] |
  \item[Yi-chku] Chin Tsu-tung 金 祖 同 . Yin-ch’i yi-chu 股 契 遺 珠 . Shanghai, 1939. Reprint, 291 Taipei, 1975. J po 易編 Tung Tso-pin. Hstao-t’un ti-erh-pen: Yin-hsti wen-tzu: 人 2-222 小 地 第 二 本 : BBG iE
  \item[CF:] ZH. Pt. 1 [Nanking], 1948. Pt. 2 [Nanking], 1949. Pt. 3, Taipei, 1953.
  \item[Yi-ts‘un] Shang Ch’eng-tso. Yin-ch’i yi-ts’un BY ARR AE. Nanking, 1933. Reprint, Tokyo, 1966. | ; Yin-hsii\_ Ming Yi-shih [ James Mellon Menzies]. Yin-hsii pu-tz’u HR Mi b BY (Oracle Records Lig from the Waste of Yin). Shanghai, 1917. Reprint, Taipei, 1972. \{ 4 1 | Lo | t
\end{description}
